\chapter{The Economics of Belief Production}
\label{ch:economics-belief}

\epigraph{There is no such thing as a free market.}{Ha-Joon Chang}

\noindent
Belief production has costs, returns, market structures, and welfare effects.
This chapter applies economic analysis to the doxonomic system itself, asking: what does belief cost, who earns the returns, and what are the welfare consequences of concentrated reality production?

\section{Cost Functions for Narrative Repetition}
\label{sec:ch10-costs}

\begin{model}[Stylized Doxonomic Cost Function]
\label{mod:cost-function}
\index{cost function}
A simple candidate cost function for achieving belief level $\mu$ for proposition $p$ through channel $k$ is
\[
  c_k(\mu) = \frac{-\ln(1-\mu)}{\credibility{k} \cdot \reach{k}}
\]
This is a heuristic, not an empirically derived function.
Its justification is qualitative: the last few percentage points of adoption are exponentially more expensive than the first, because holdouts are more resistant, more isolated, or more committed to alternatives.
The functional form is chosen for tractability.
Empirical estimation of actual cost curves is an open problem for the field.
\end{model}

The total cost of achieving population-level belief $\mu(p) = \bar{\mu}$ using $m$ channels is the solution to
\[
  \min_{\mu_1, \ldots, \mu_m} \sum_{k=1}^m c_k(\mu_k) \quad \text{subject to} \quad \sum_k w_k \mu_k = \bar{\mu}
\]
where $w_k$ is the population weight of group $k$.

\section{Ideological Return on Investment}
\label{sec:ch10-iroi}

\begin{definition}[Ideological ROI]
\label{def:iroi}
\index{ideological ROI}
The \emph{ideological return on investment} for belief $p$ through channel $k$ is
\[
  \iroi_k(p) = \frac{\Delta \belief{p}}{\Delta \funding{k}} = \frac{\credibility{k} \cdot \relevance{k}{p} \cdot \reach{k} \cdot \persistence{k}}{c_k'(\mu)}
\]
evaluated at the current belief level.
The aggregate $\iroi$ is the weighted average across all active channels.
\end{definition}

\begin{proposition}[Decreasing IROI]
\label{prop:decreasing-iroi}
$\iroi_k(p)$ is decreasing in $\mu(p)$: beliefs that are already widely held offer lower marginal returns to additional investment.
This creates an incentive to diversify across beliefs or to shift from persuasion to maintenance once capture is achieved.
\end{proposition}

\section{Belief as Public Good or Public Bad}
\label{sec:ch10-public-good}

Beliefs are non-excludable (you cannot prevent someone from adopting a freely circulating narrative) and partially non-rival (one person's belief does not deplete the narrative).
This gives beliefs the structure of a \emph{public good}---or, when harmful, a \emph{public bad}.

\begin{definition}[Doxonomic Externality]
\label{def:doxonomic-externality}
\index{doxonomic externality}
A \emph{doxonomic externality} arises when a subsidized belief $p$ imposes costs on agents who did not choose to produce or receive it.
Examples include: reduced social trust from selfishness narratives, increased tolerance for inequality from meritocratic myths, or reduced collective action capacity from individualist anthropologies.
\end{definition}

\begin{theorem}[Welfare Loss from Doxonomic Externalities]
\label{thm:welfare-loss}
If the social cost of belief $p$ is $D(\mu(p))$ (increasing and convex in adoption level) and private funders do not internalize this cost, the equilibrium adoption level $\mu^*$ exceeds the social optimum $\mu^{**}$:
\[
  \mu^* > \mu^{**} \quad \text{where} \quad \mu^{**} = \arg\min_{\mu} \left[ c(\mu) + D(\mu) \right]
\]
The welfare loss is $\int_{\mu^{**}}^{\mu^*} [D'(\mu) - c'(\mu)] \, d\mu > 0$.
\end{theorem}

\section{Epistemic Herfindahl Index}
\label{sec:ch10-ehhi}

\begin{definition}[EHI]
\label{def:ehhi}
\index{epistemic Herfindahl index!formal}
The \emph{epistemic Herfindahl index} is
\[
  \ehhi = \sum_{k=1}^{m} s_k^2 \qquad \text{where} \qquad s_k = \frac{\Theta_k}{\sum_{j} \Theta_j}
\]
and $\Theta_k = \funding{k} \cdot \credibility{k} \cdot \reach{k}$ is the effective doxonomic throughput of channel $k$.
\end{definition}

High $\ehhi$ indicates a concentrated market for reality production, in which a small number of institutions control most of the population's narrative exposure.
In standard antitrust economics, $\mathrm{HHI} > 0.25$ signals high concentration.
Whether the same threshold applies to belief production is an empirical question---narrative markets differ from commodity markets in credibility dynamics, network effects, and audience resistance.
We adopt $\ehhi > 0.25$ as a provisional warning indicator, not as a settled regulatory standard.

\section{Subsidy Structures}
\label{sec:ch10-subsidies}

Much doxonomic spending is invisible because it is structured as tax-exempt philanthropy, corporate sponsorship, or institutional cross-subsidy.

\begin{model}[Hidden Doxonomic Transfer]
\label{mod:hidden-transfer}
Let a corporation spend $S$ on think-tank funding.
The effective doxonomic transfer is $S \cdot (1 - t)$ where $t$ is the tax rate on the deduction.
Society effectively co-finances the doxonomic investment through foregone tax revenue.
The net public cost is $S \cdot t$ for a narrative chosen by the funder.
\end{model}

\section{Notes and References}
\label{sec:ch10-notes}

The economics of public goods and externalities draws on \textcite{stiglitz2012} and \textcite{arrow1963}.
The market for ideas framework is from \textcite{coase1974}.
Media market concentration is analyzed in \textcite{bagdikian2004}.
The Herfindahl index originates in \textcite{herfindahl1950}.
On manipulation as an economic phenomenon, see \textcite{akerlof2015}.
The political economy of think-tank funding is documented in \textcite{mayer2016}.

\section{Exercises}

\begin{exercise}
Compute $\ehhi$ for the following hypothetical market: three channels with throughputs $\Theta_1 = 50$, $\Theta_2 = 30$, $\Theta_3 = 20$.
Is this market concentrated by doxonomic standards?
\end{exercise}

\begin{exercise}
A funder can invest \$10M in advertising ($\credibility{} = 0.2$, $\reach{} = 0.6$) or in university research ($\credibility{} = 0.85$, $\reach{} = 0.05$).
Compute $\iroi$ for each.
Under what conditions is the university investment more effective despite lower reach?
\end{exercise}

\begin{exercise}
Model the welfare loss from the belief ``poverty reflects personal failure'' using a specific damage function $D(\mu) = \alpha \mu^2$.
Find the socially optimal belief level and the deadweight loss.
\end{exercise}

\begin{exercise}
Estimate the hidden public subsidy to ideological production through the U.S.\ 501(c)(3) tax exemption.
What total doxonomic spending does this structure enable?
\end{exercise}
